\section{Continous Integration and development process}
Based on pre-study and requirements choices an starting architecture and pipeline was designed. For the architecture choice a web app was chosen due to its accessibility and scalability.
To support development of the platform developing a pipeline to give feedback and quality. 
At the start it was prioritized to have a pipeline that gives us automated feedback on our implementations.
Technical details about the platform and pipeline and discussion about how they impacted the project. 

\subsection{CI Pipeline Overview}
With the project work on GitHub it was considered using GitHub Actions to design the pipeline.
As the team was working on modern IDE which did small checks before pushing like static code analysis.
Building the application was also done frequently locally. 
As a lot of checks was done before pushing, but in the cases where it was not a pipeline was create to catch these errors.
In the hopes of catching unseen errors and issues we designed the pipeline as follows: 
The skeleton for the pipeline was from gradle and these steps were added:
\begin{itemize}
  \item On each pull to repository
  \item Build using Gradle
  \item Run tests using JUnit
\end{itemize}
additionally, it will provide feedback linked to each commit. By being connceted to each commit and getting feedback, we can proactiviely address any issues that arise during the build and test process.
We plan to add more steps as testing and integration evolves. additionally as development progresses towards a release, we need reliable generation of artifacts that have passed the pipeline and ready for release. 
Currently, we create artifacts manually for release of now we are doing creation of artifacts manually on each release.

% later version
%\subsection{placeholder The pipeline}
%For the CI pipeline, we chose GitHub Actions.
%This was a practical choice because GitHub Actions is both familiar and widely used, and our project is already hosted on GitHub.
%As a result, the pipeline was relatively easy to set up and integrates naturally with the rest of our workflow.

%At the current stage, the main purpose of the pipeline is to automate building and testing whenever code is pushed.
%This allows us to detect problems earlier and gives us better insight into when a particular commit may have introduced an error.
%Since multiple team members are working on the project simultaneously, this is important for debugging and coordination.
%Without automated feedback, it can become difficult to identify where a problem originated.

%The pipeline is also useful for release preparation.
%Running builds and tests before or during merges gives a faster indication of whether the code base is
%stable enough to move toward a release.
%At a later stage, the pipeline can be extended with additional features such as test coverage reporting and artifact generation.
%When the frontend and backend are more mature, deployable artifacts can also be produced so that others can
%test the system more easily and releases can be tracked more systematically. 

\section{System Architecture} \subsection{HighLevel Design}

The system consists of a frontend application and a backend service that communicates with the frontend
via HTTP/JSON. The backend exposes RESTful APIs which the frontend makes requests to in order to retrieve data.
A relational database (PostgreSQL) is used to store all persistent

\subsection{Choice of Languages and Frameworks}

For backend development, we chose Kotlin.
One reason for this decision is that the group is already familiar with Java,
and Kotlin works within the same ecosystem.
Since Spring is a very common framework and Kotlin integrates well with it, this combination was a natural choice.
Kotlin also provides a more modern syntax while retaining access to mature Java-based libraries and tools.
% use latex tools to show stack in a list or diagram
This is our entire stack.
\begin{itemize}
    \item Kotlin for the programming language
    \item Spring Framework for backend development
    \item PostgreSQL for databases
    \item Docker for container management and orchestration
    \item React and TypeScript for frontend development
\end{itemize}
The Spring Framework has multiple additions like security, web server and JPA for ORM that makes it a good choice for developing a REST api for our frontend.

Docker is essential for ensuring that our application runs in an isolated environment, 
regardless of the underlying infrastructure. This makes it easy to 
To ensure the
For backend testing and web services, we use Spring Web together with JUnit.
This provides a familiar testing environment and supports the development of backend functionality in a structured way.

For the frontend, we use React with TypeScript.
React was selected because it supports component-based design and is widely used in web application development.
TypeScript helps improve code quality and maintainability through stronger type checking.
We also use React Router in order to structure the application into pages and make navigation easier to manage.

At the current stage, the frontend and backend are still being developed somewhat independently in their own development environments.
This means that some integration remains unfinished, and in some cases mocking or partial functionality
may be necessary until the systems are connected more fully.

\subsection{Git Workflow and Version Control}

We use Git for version control so that work can be saved, tracked, edited, and reversed when necessary.
Since both frontend and backend are developed in parallel, version control is essential for managing changes
across the system and coordinating team contributions.

Our current workflow uses a main branch for stable and project-wide changes,
such as CI updates and release-related work.
In addition, we maintain separate frontend and backend branches where ongoing development in each area can continue more independently.
For larger pieces of functionality, smaller branches may also be created from these development branches.

This workflow makes it possible to work on different features for frontend or backend or both at the same time. 
Merging becomes necessary when functionality spans both parts of the system and immediately before releases.
At a later stage, this workflow should be strengthened further with clearer tagging to each feature and version tagging.

\subsection{Implemented User Stories and Current Functionality}

The current development work is based on user stories defined from the perspectives of students and lecturers.
These stories reflect the main goals of the project, which are to improve student engagement, support participation,
and create a more motivating study experience.

The user stories currently guiding the project include:

\begin{itemize}[leftmargin=1.5em]
    \item As a Student, I want to see my current level and progress in a course without relying on grades.
    \item As a Student, I want to ask anonymous questions without feeling stupid.
    \item As a Lecturer, I want to see the students engage in the subject by asking questions.
    \item As a Student, I want streaks for weekly engagement so I do not lose consistency.
    \item As a Student, I want to participate in study groups so I can interact with other students.
    \item As a Student, I want to find study groups with similar goals.
    \item As a Student, I want to receive feedback if my participation drops.
    \item As a shy student, I want the option to compete anonymously so that I feel safe.
    \item As a distracted student, I want a timed focus mode so that I avoid interruptions.
    \item As a motivated student, I want to maintain a study streak so that I build consistent habits.
\end{itemize}

These user stories show that the project is not only about academic administration, but also about motivation, engagement, and participation. Some stories focus on individual self-regulation, such as progress tracking, streaks, and focus mode. Others focus on social and interactive aspects, such as study groups, anonymous questions, and lecturer insight into engagement.

At this stage, not all of these stories have been fully implemented.
Some are still being developed, and some may currently exist only as partial functionality, prototype logic, or planned features.
However, they provide a clear structure for prioritization and sprint work.

\subsection{CI Pipeline Reflection}
A lot of short-lived feature branches was made during development and the pipeline was doing everything on each commit a change was done.
Because our budget had to be controlled and the short-lived feature branches was tested locally.
A pipeline that ran builds, artifacts and tests only in main development branches
(frontend, main, backend) to get feedback on pull requests and have working artifacts.
The short-lived branches and frequent merges into main development branches with a pipeline gave valuable feedback on features and ensured quality with testing.

The Taming googles large scale CI paper had a somewhat same problem as they needed to reduce time
on testing and rather only run tests which is prone to find faults rather than running tests that always come out true. \cite{GoogleTesting}
For our case we cut down on CI runs and put more effort in that developers should thoroughly
check their code on their feature. Google also made developer do extra checks when affected by faulty code or high risk code. \cite{GoogleTesting}

React testing as automating it and doing integration testing seems difficult to do using the pipeline.
The development is right now at a small scale, so making integration testing has not been prioritized as getting
a MVP is number one priority, but at a larger scale this doing as much automated testing as possible is beneficial
for quicker releases and feedback.

Google's large scale CI pipeline started from a small scale thinking and was able to make adequate changes
to the pipeline and version control, so it could handle the scalability problems it faced.
At the start they had no idea of how massively it would scale. Our pipeline went through optimization as the building stage did not have to run at every stage.
Also, putting the developer more responsible of their code as modern IDEs run typical checks before commiting, so developers are more aware for what could cause a breakage.
The approach we had with version control worked well for us. Google shows that it is possible to contain everything within one repository, so the version control has possibilites to improve for scalability. \cite{GoogleRepo}
The pipeline has to follow how the version control is done and both systems works together. 
A plan for future development is to make both systems ready for possible scaling, so large change in development would not destroy the repository.

When scaling our pipeline to enterprise levels a lot of problems would occur.
Our testing only has automated unit testing and everything else is done manually. 
The pipeline will only do unit testing and there would be a lot of manual testing. 
By making integration testing and some automated release testing based on features, so user story based testing would work well with the pipeline.
As it would add more dependencies to building the application, manual building would be too time consuming for production. And the pipeline would need probably run building in parallel on different branches to handle different features and remove unused dependencies.

